\chapter{Introduction}

Quadratic integer sequences play a central role in number theory,
particularly in the study of congruence properties, polygonal numbers,
and prime number distributions. The present work considers the special
quadratic sequence
\begin{equation}\label{eq:k_main}
k(n) = \frac{1}{6}(2n^2 + 3n + 1),
\end{equation}
which possesses a Diophantine property: it takes integer
values if and only if $n$ lies in the congruence classes
\begin{equation}\label{eq:congruence}
n \equiv 1, 5 \pmod{6}.
\end{equation}

These two classes are of fundamental number-theoretic significance, as they
contain all prime numbers greater than~$3$ (Theorem~\ref{satz:primzahlen} in
Chapter~\ref{sec:arithmetik}).

\section*{Motivation and Objectives}

The geometric visualisation of number-theoretic structures has a long
tradition and has led to significant insights. Well-known examples include:
\begin{itemize}[leftmargin=2em]
\item The \emph{Ulam spiral}~\cite{Stein1964}
\item The \emph{Sacks spiral}~\cite{Sacks2003}
\item Quadratic forms in prime number theory
\end{itemize}

While the Ulam and Sacks spirals are primarily empirically motivated ---
patterns observed first, explanations sought afterwards --- the theory of
quadratic forms represents the complementary algebraic tradition; indeed,
the diagonal patterns of the former are governed by quadratic polynomials.
The present work pursues this second direction:
starting from the quadratic sequence~\eqref{eq:k_main}, a geometric
structure is \emph{derived}, not merely visualised. The contribution is accordingly a geometric synthesis of known arithmetic building blocks along a single derivation path, rather than new number theory.

\clearpage

\section*{Main Results}

The central results of this work are:
\begin{enumerate}
\item Alternating minor and major steps with closed-form formulae
\item Circle circumferences with constant radius difference $\Delta R = \frac{12}{\pi}$
\item Archimedean spiral $r(\theta) = \frac{6}{\pi^2}\theta$
\item Parametrisation via $p$: $r(p)=\frac{6}{\pi}\sqrt{p/3}$, $\theta(p)=\pi\sqrt{p/3}$
\item Exact arc-length integral
\item Conical helix as a 3D embedding
\item All primes $p>3$ appear as a distinguished discrete subset of the helix (corollary of the congruence class structure)
\end{enumerate}

\section*{Structure of the Work}

The work is organised into five parts, mirroring the derivation path.
\textbf{Part~I}, \emph{The Parabola}
(Chapters~\ref{sec:arithmetik}--\ref{sec:distanz_sequenzen}), develops
the arithmetic foundations: congruence condition, the closed-form step
formulae, and Euclidean distances. \textbf{Part~II}, \emph{The
Archimedean Spiral} (Chapters~\ref{sec:spirale}--\ref{ch:taylor}),
derives the spiral from the minimality of the pairing and treats
parametrisation, arc length, and the Taylor approximation.
\textbf{Part~III}, \emph{The Cone} (Chapters~\ref{ch:cone}
and~\ref{sec:kegelstrahlen}), constructs the enveloping cone and its ray
geometry. \textbf{Part~IV}, \emph{The Conical Helix}
(Chapters~\ref{ch:helix}--\ref{ch:examples}), assembles the helix,
derives the prime number corollary, and verifies the arc lengths.
\textbf{Part~V}, \emph{The Direct Solution}
(Chapter~\ref{sec:visualisierung} onwards), presents the exact
$Q$-parametrisation, the direct parabola--helix construction
(Chapter~\ref{ch:parabola_helix}), the summary
(Chapter~\ref{sec:zusammenfassung}), and the formula collection.

\section*{Notation}

\begin{itemize}[leftmargin=2em]
\item $k(n)$: quadratic mean-value sequence
\item $\delta_k$: difference sequence of the integer values
\item $\Delta_{\mathrm{Major}},\ \Delta_{\mathrm{Minor}}$: major and minor step sizes
\item $a_k$: closed-form generator of $(\delta_k)$
\item $Q(k) = \sqrt{48k+1}$: universal quantity
\item $a = 6/\pi^2$: spiral slope; $c_1 = 3/\pi$: height rate
\item $\theta$: angular parameter
\item $p$: function-value parameter
\item $r(\theta)$: radius function
\item $\gamma(\theta)$: helix parametrisation
\end{itemize}
